% !TEX program = pdflatex
\documentclass[11pt]{article}

% ==== Packages ====
\usepackage[utf8]{inputenc}
\usepackage[T1]{fontenc}
\usepackage{geometry}
\geometry{margin=1in}
\usepackage{microtype}
\usepackage{amsmath,amssymb,amsthm,mathtools}
\usepackage{bm,bbm}
\usepackage{enumitem}
\usepackage{hyperref}
\usepackage[nameinlink,capitalize,noabbrev]{cleveref}
\usepackage{xcolor}
\usepackage{algorithm}
\usepackage{algpseudocode}
\usepackage{listings}
\lstset{basicstyle=\ttfamily\small,breaklines=true,frame=single,language=Python}

% ==== Theorem Environments ====
\theoremstyle{definition}
\newtheorem{definition}{Definition}[section]
\newtheorem{example}[definition]{Example}
\theoremstyle{plain}
\newtheorem{theorem}[definition]{Theorem}
\newtheorem{proposition}[definition]{Proposition}
\newtheorem{lemma}[definition]{Lemma}
\newtheorem{corollary}[definition]{Corollary}
\theoremstyle{remark}
\newtheorem{remark}[definition]{Remark}

% ==== Macros ====
\newcommand{\Sig}{\mathrm{Sig}}
\newcommand{\A}{\mathcal{A}}
\newcommand{\Hcal}{\mathcal{H}}
\newcommand{\Pbb}{\mathbb{P}}
\newcommand{\Qbb}{\mathbb{Q}}
\newcommand{\Zbb}{\mathbb{Z}}
\newcommand{\Cbb}{\mathbb{C}}
\newcommand{\Mfun}{M_{\mathrm{fun}}}
\newcommand{\Mint}{M_{\mathrm{int}}}
\newcommand{\Mhat}{\widehat{M}}
\DeclareMathOperator{\Tr}{Tr}
\DeclarePairedDelimiter{\norm}{\lVert}{\rVert}
\DeclarePairedDelimiter{\abs}{\lvert}{\rvert}

\begin{document}
\begin{titlepage}
    \centering

    {\Huge\bfseries Prime-Graded Multiplicity}\\
    \vspace{0.5cm}
    {\LARGE\bfseries Unified Specification and Formalization \par}

    \vspace{2cm}
    {\LARGE\itshape Ryan O. Van Gelder \par}
    \vspace{1cm}
    {\Large Citizen Gardens} \\
    \vspace{0.5cm}
    \Large{Institute for Mathematical Discovery} \\ 
    \vspace{0.5cm}
    {\normalsize \today \par}

 \vfill

    {\bfseries Abstract \par}
    \vspace{0.5em}
    \begin{minipage}{0.85\textwidth}
        \small
We formalize the Prime\,–\,Graded Multiplicity (PGM) for the multiplicity operator $\mathbf M$ as a single categorical object with consistent algebraic, analytic, and representation\,–\,theoretic faces. The core consists of a signature category $\Sig$, a prime\,–\,graded $*$\,–\,algebra $(\A,\{P_p\})$ with a positive state $\varphi$, and a functorial signature extractor $F$. From this backbone we derive: (i) the strict monoidal functor $\Mfun: \Sig\to \Qbb^{\times}$; (ii) the integer ledger $\Mint$ and the state\,–\,dependent sector distribution $\Mhat(T,p)$; (iii) the diagonal representation $\mathbf M$ acting on a signature Hilbert space with number operators $N_p$. We prove normalization, stability, and equivalence properties, give algorithms and complexity, and delimit a multiple\,–\,frame extension as future work.
 \end{minipage}
\end{titlepage}

\newpage

\tableofcontents



\section{Introduction}
The multiplicity operator $\mathbf M$ has appeared informally in several contexts. This article provides a rigorous, \emph{unified} specification in which $\mathbf M$ is not many different objects but one categorical functor whose other appearances are canonical incarnations. The resulting scaffold separates exact algebraic information (integer ledgers) from analytic, state\,–\,dependent summaries (normalized sector distributions) and from diagonal operator representations. Speculative extensions (e.g. non\,–\,abelian behavior) are quarantined and given precise mathematical triggers (frame misalignment).

\paragraph{Contributions.} (1) A strict monoidal functor $\Mfun$ on signature data and its valuation properties; (2) a computable, basis\,–\,independent definition of $\Mhat(T,p)$ via projectors and a state; (3) a diagonal representation with $\log \mathbf M=\sum_p (\log p)N_p$; (4) stability and convergence criteria; (5) algorithms and worked examples.

\section{Preliminaries and Notation}
We write $\Pbb$ for the set of primes; $\Zbb,\Qbb,\Cbb$ have their usual meanings. All sums over $p$ are over primes; supports are finite unless stated otherwise. For $X$, $\norm{X}$ denotes an operator or Frobenius norm depending on context; $\Tr$ denotes the trace when $\A\subseteq \Cbb^{n\times n}$.

\section{Backbone of the PGM}
\subsection{Signature Category $\Sig$}
\begin{definition}[Signature category]
Let $\Sig=\{e:\Pbb\to\Zbb\text{ with finite support}\}$. Equip $\Sig$ with $(e\otimes f)(p)=e(p)+f(p)$, unit $0$, and dual $e^{\vee}=-e$. Then $(\Sig,\otimes,0)$ is a strict commutative monoid, viewable as a strict monoidal category with one object.
\end{definition}

\subsection{Prime\,–\,Graded $*$\,–\,Algebra with State}
\begin{definition}[Prime grading and state]
Let $\A=\bigoplus_{p\in\Pbb}\A_p$ be a complex $*$\,–\,algebra with orthogonal projectors $P_p:\A\to\A_p$ satisfying $P_pP_q=\delta_{pq}P_p$ and $\sum_p P_p=I$ (finite sum on supports). Let $\varphi: \A\to\Cbb$ be positive and normalized.
\end{definition}

\subsection{Signature Extractor}
\begin{definition}[Functorial extractor]\label{def:extractor}
A map $F:\operatorname{Obj}(\A)\to\Sig$ satisfies
\begin{equation}\label{eq:Faxioms}
F(XY)=F(X)+F(Y),\qquad F(X^{*})=-F(X),\qquad F(\mathbf 1)=0.
\end{equation}
When restricted to multiplicatively factorable objects, $F$ records prime\,–\,exponent ledgers.
\end{definition}

\section{The $\mathbf M$ Object and Its Faces}
\subsection{Functorial $\Mfun$ (Categorical Backbone)}
\begin{definition}[Functorial $\Mfun$]\label{def:Mfun}
Define $\Mfun: \Sig\to \Qbb^{\times}$ by $\Mfun(e)=\prod_p p^{e(p)}$. Then $\Mfun$ is strict monoidal: $\Mfun(0)=1$, $\Mfun(e+f)=\Mfun(e)\Mfun(f)$, $\Mfun(-e)=\Mfun(e)^{-1}$.
\end{definition}
\begin{proposition}[Valuation property]\label{prop:valuation}
Let $\nu_p: \Qbb^{\times}\to\Zbb$ be the $p$\,–\,adic valuation. Then $\nu_p\circ \Mfun=\pi_p: \Sig\to\Zbb$, where $\pi_p(e)=e(p)$.
\end{proposition}
\begin{proof} Immediate from the multiplicative definition of $\Mfun$ and the uniqueness of prime factorization in $\Qbb^{\times}$.\end{proof}
\begin{remark}[Exponential form]
Acting diagonally on signature labels, $\Mfun=\exp\big(\sum_p (\log p)\,\pi_p\big)$.
\end{remark}

\subsection{Integer Ledger and State\,–\,Dependent Distribution}
\begin{definition}[Integer ledger and distribution]\label{def:ledger-dist}
Given $X$ with $e=F(X)$, define the integer ledger $\Mint(X,p):=e(p)\in\Zbb$. For $T\in\A$, define the sector mass and normalized weight
\begin{equation}\label{eq:sector-weight}
\mu_p(T):=\varphi\big((P_pT)^{*}(P_pT)\big)\ge 0,\qquad \Mhat(T,p):=\frac{\mu_p(T)}{\sum_q\mu_q(T)}\in[0,1].
\end{equation}
\end{definition}
\begin{proposition}[Normalization]\label{prop:normalization}
If $T\neq 0$, then $\sum_p \Mhat(T,p)=1$.
\end{proposition}
\begin{proof} Positivity and linearity of $\varphi$ plus orthogonality $P_pP_q=\delta_{pq}P_p$ yield $\sum_p\mu_p(T)=\varphi\big((\sum_p P_pT)^*(\sum_p P_pT)\big)=\varphi(T^*T)>0$.\end{proof}
\begin{remark}[Ledger vs. distribution]
$\Mint$ is exact, algebraic, and functorial; $\Mhat$ is analytic and state\,–\,dependent. They agree on pure signatures up to a chosen normalization rule (e.g. proportional to $\abs{e(p)}$) but generally differ on superpositions.
\end{remark}

\subsection{Diagonal Representation on Signature Space}
\begin{definition}[Signature Hilbert space and number operators]
Let $\Hcal=\mathrm{span}\{\,|e\rangle: e\in\Sig\,\}$ and define number operators $N_p|e\rangle=e(p)|e\rangle$. Define $\mathbf M$ by
\begin{equation}\label{eq:diagM}
\mathbf M\,|e\rangle=\Mfun(e)\,|e\rangle,\qquad \log\mathbf M=\sum_p (\log p)\,N_p.
\end{equation}
\end{definition}
\begin{proposition}[Monoidal–Diagonal Equivalence]\label{prop:gauge}
Let $\rho: \Sig\to\mathcal U(\Hcal)$ be the regular diagonal representation $\rho(e)|f\rangle=|e+f\rangle$. Then $\mathbf M=\sum_e \Mfun(e)|e\rangle\langle e|$ is the unique diagonal operator (up to a central scalar) whose conjugation $X\mapsto \mathbf M X \mathbf M^{-1}$ implements a signature\,–\,scale gauge compatible with \cref{def:Mfun}.
\end{proposition}
\begin{proof} Diagonality and \cref{def:Mfun} force the action on the eigenbasis $|e\rangle$; uniqueness follows from the spectral theorem for commuting $\{N_p\}$.\end{proof}

\section{Stability and Convergence}
\begin{definition}[Weighted prime response]
For real parameters $\{\alpha_p\}$ define
\begin{equation}\label{eq:Salpha}
S_{\alpha}(T):=\sum_p \Mhat(T,p)\,p^{\alpha_p}.
\end{equation}
\end{definition}
\begin{proposition}[Bounds and convergence]\label{prop:bounds}
If only finitely many primes are active in $T$, then $0\le S_{\alpha}(T)\le \max_p p^{\alpha_p}$. If infinitely many are active, convergence of \cref{eq:Salpha} requires $\alpha_p\to-\infty$ sufficiently fast; otherwise $S_{\alpha}(T)$ is undefined.
\end{proposition}
\begin{proposition}[Contractivity gate]\label{prop:contract}
In regimes where an update $T\mapsto \mathcal U(T)$ linearizes with gain aligned to the weights in \cref{eq:Salpha}, a sufficient stability condition is $\abs{S_{\alpha}(T)}<1$.
\end{proposition}

\section{Computation and Algorithms}
\begin{algorithm}[H]
\caption{Normalized sector weights $\Mhat(T,p)$}
\label{alg:weights}
\begin{algorithmic}[1]
\Require projectors $\{P_p\}$, state $\varphi$, element $T\in\A$ with finite prime support
\ForAll{active primes $p$}
  \State $T_p\gets P_p T$
  \State $\mu_p\gets \varphi\big(T_p^{*}T_p\big)$
\EndFor
\State $\mu\gets \sum_q \mu_q$; return $\Mhat(T,p)=\mu_p/\max(\mu,\varepsilon)$ for each $p$ (small $\varepsilon>0$)
\end{algorithmic}
\end{algorithm}
\noindent\textbf{Complexity.} Linear in the number of active primes. If projectors are approximate, report an orthogonality defect $\delta=\sum_{p\ne q}\norm{P_pP_q}$ and propagate its effect on $\Mhat$.

\section{Multiple\,–\,Frame Extension (Future Work)}
Let $\{P_p^{(i)}\}$ be alternative prime frames with intertwiners $U_{ij}$ satisfying $P_p^{(j)}=U_{ij}P_p^{(i)}U_{ij}^{-1}$. Define
\begin{equation}
\Mhat^{(i)}(T,p)=\frac{\varphi\big((P_p^{(i)}T)^{*}(P_p^{(i)}T)\big)}{\sum_q\varphi\big((P_q^{(i)}T)^{*}(P_q^{(i)}T)\big)},\qquad \mathbf M^{(i)}=\exp\Big(\sum_p (\log p)\,N_p^{(i)}\Big).
\end{equation}
Then $\mathbf M^{(j)}=U_{ij}\,\mathbf M^{(i)}\,U_{ij}^{-1}$ and $[\mathbf M^{(i)},\mathbf M^{(j)}]=0$ iff the frames are simultaneously diagonalizable; non\,–\,commutation measures frame misalignment.

\section{Worked Examples}
\begin{example}[Arithmetic toy model]
Let $e(2)=1,\ e(3)=-1,\ e(5)=2$. Then $\Mint(X,2)=1$, $\Mint(X,3)=-1$, $\Mint(X,5)=2$ and $\Mfun(F(X))=2^{1}3^{-1}5^{2}=50/3$. Choosing $\Mhat$ from absolute exponents gives $\Mhat(X,2)=1/4$, $\Mhat(X,3)=1/4$, $\Mhat(X,5)=1/2$.
\end{example}
\begin{example}[Matrix model with trace state]
Let $\A\subseteq \Cbb^{n\times n}$, $\varphi=\Tr$, and $P_2,P_3,P_5$ be block projectors. For $T=\mathrm{diag}(A,B,C)$ we have $\mu_2=\norm{A}_F^2$, $\mu_3=\norm{B}_F^2$, $\mu_5=\norm{C}_F^2$ and $\Mhat(T,2)=\norm{A}_F^2/(\norm{A}_F^2+\norm{B}_F^2+\norm{C}_F^2)$.
\end{example}

\section{Deployment Checklist}
Specify: (1) the graded algebra $\A$ and state $\varphi$; (2) projectors $\{P_p\}$ and an orthogonality\,–\,defect budget; (3) a functorial extractor $F$ satisfying \cref{def:extractor}; (4) a rule linking ledgers to distributions on pure signatures (if desired); (5) exponents $\{\alpha_p\}$ and convergence guarantees for $S_{\alpha}$.

\appendix
\section{Reference Implementation (Python)}
\begin{lstlisting}
from dataclasses import dataclass
from typing import Dict, Callable, Iterable
import numpy as np

Scalar = float | complex
Projector = Callable[[np.ndarray], np.ndarray]
State = Callable[[np.ndarray], Scalar]
Signature = Dict[int, int]  # prime -> exponent

@dataclass
class PrimeGradedFramework:
    primes: Iterable[int]
    projectors: Dict[int, Projector]  # p -> P_p
    state: State                      # φ: A -> C (e.g., np.trace)

    @staticmethod
    def M_fun(signature: Signature) -> float:
        out = 1.0
        for p, exp in signature.items():
            out *= (p ** exp)
        return out

    def sector_weight(self, T: np.ndarray, p: int) -> float:
        T_p = self.projectors[p](T)
        mu_p = float(self.state(T_p.conj().T @ T_p))
        total = 0.0
        for q in self.primes:
            T_q = self.projectors[q](T)
            total += float(self.state(T_q.conj().T @ T_q))
        eps = 1e-12
        return mu_p / max(total, eps)

    def S_alpha(self, T: np.ndarray, alpha: Dict[int, float]) -> float:
        return sum(self.sector_weight(T, p) * (p ** alpha[p]) for p in self.primes)
\end{lstlisting}

\section{Lipschitz Continuity: A Proof Sketch}
Assuming $\varphi$ uniformly continuous and $\norm{P_p}\le 1$, expand $\Mhat(T,p)$ using \cref{eq:sector-weight} and apply standard quotient stability bounds to obtain
\[
\abs{\Mhat(T,p)-\Mhat(S,p)}\ \le\ \frac{C\,\norm{T-S}}{\min(\norm{T}^2,\norm{S}^2)}
\]
for some $C$ depending on $\varphi$.

\section{Scope and Non\,–\,Claims}
This formalization is purely mathematical; any physical modeling (e.g. dissipative/ conservative splittings) must be justified independently from a chosen dynamical principle (Lyapunov, unitarity, etc.).

\appendix
\section{A Free Abelian Structure of Signatures on Primes}

We develop a certified mathematical framework where prime factorization provides a canonical encoding of tensor structure. Objects are \emph{prime signatures} --- finitely supported integer labelings of the set of primes --- equipped with a strict symmetric monoidal structure (tensor = pointwise addition, unit = $0$, dual = negation). A central \emph{multiplicity functor} $\M:\Sig\to \Q^{\times}$ maps signatures to units of the rationals by $\M(e)=\prod_{p} p^{e_p}$, thereby transporting tensor to multiplication and duals to inversion. We formalize the key laws ($\M(0)=1$, $\M(e+f)=\M(e)\,\M(f)$, $\M(-e)=\M(e)^{-1}$) and show prime-conservation for morphisms in the discrete categorical model. This yields algebraic certificates for tensor product and contraction that are independent of floating-point numerics and integrate seamlessly with an ACE control loop for stability-aware contraction planning.

\section{Introduction}
\label{sec:intro}
Prime factorization is canonical, global, and choice-free. We leverage this to encode the structural ``type'' of tensor axes and to define a multiplicative invariant that certifies the correctness of tensor operations. The resulting framework---\emph{Typed-Prime-Signatures} (TPS)---unifies:
\begin{itemize}[leftmargin=*]
  \item a free abelian structure of signatures on primes,
  \item a monoidal/dual semantics (tensor/dual $\leftrightarrow$ add/negate exponents),
  \item a functorial multiplicity invariant into $\Q^{\times}$ (group of rational units), and
  \item certified computation: tensor product and contraction preserve the invariant and hence cannot create or destroy ``prime mass".
\end{itemize}
We also outline an ACE-integrated control loop that uses these algebraic guarantees as hard constraints while learning numerically stable contraction orderings.

\paragraph{Contributions.}
(1) A strict symmetric monoidal skeleton on prime signatures with duals; (2) a monoidal functor $\M:\Sig\to\Q^{\times}$ proving conservation laws; (3) a computational semantics that lifts signatures to axis-typed tensors; (4) certified operations and property-based testing; (5) an integration plan with ACE for stability-aware, safety-preserving tensor contraction.

\section{Prime Signatures and Monoidal Structure}
\label{sec:signatures}
Let $\PP$ denote the set of natural primes. We write $e=(e_p)_{p\in\PP}\in \Z^{(\PP)}$ for a finitely supported integer labeling of primes.

\begin{definition}[Prime signatures]
The set of \emph{prime signatures} is
\[
  \Sig \;:=\; \{ e:\PP\to\Z \;:\; e \text{ has finite support}\}\;\cong\;\Z^{(\PP)}.
\]
The support is $\supp(e):=\{p\in\PP: e_p\neq 0\}$.
\end{definition}

\begin{definition}[Monoidal and dual structure]
Define tensor by pointwise addition $(e\otimes f)_p := e_p+f_p$, the unit by $0$, and the dual by $e^\vee := -e$. Thus $(\Sig,\otimes,0,(\cdot)^\vee)$ is a strict symmetric monoidal skeleton with duals.
\end{definition}

\section{Multiplicity Functor to $\Q^{\times}$}
\label{sec:multiplicity}
\begin{definition}[Multiplicity]
Define $\M:\Sig\to \Q^{\times}$ by
\[
  \M(e)\;:=\; \prod_{p\in\supp(e)} p^{\,e_p}.
\]
Because $\supp(e)$ is finite, this product is well-defined. Negative exponents are interpreted via $p^{e_p}=1/p^{-e_p}$.
\end{definition}

\begin{theorem}[Monoidality and duality]
\label{thm:monoidal}
For all $e,f\in\Sig$:
\[
  \M(0)=1,\qquad \M(e+f)=\M(e)\,\M(f),\qquad \M(-e)=\M(e)^{-1}.
\]
\end{theorem}
\begin{proof}
Immediate from unique factorization and the laws of exponents. Finite support ensures all products are finite.
\end{proof}

\begin{remark}[Valuations]
The coordinate $e_p$ coincides with the $p$-adic valuation $v_p(\M(e))$. Thus $e\mapsto \M(e)$ packages the family of valuations $(v_p)_p$ into a single multiplicative invariant.
\end{remark}

\paragraph{Unsigned size (optional).} Define $\Mabs(e):=\prod_p p^{\lvert e_p\rvert}\in \Q_{\ge 0}$. Then $\Mabs(e+f)=\Mabs(e)\,\Mabs(f)$, but $\Mabs(-e)=\Mabs(e)$, so $\Mabs$ ignores variance (duals) and is not a duality-preserving functor.

\section{Categorical Semantics and Conservation}
\label{sec:category}
To obtain lightweight, fully certified laws, we view signatures as objects of the \emph{discrete} category $\D(\Sig)$: morphisms are equalities.

\begin{proposition}[Prime conservation in the discrete model]
Any isomorphism $e\cong f$ in $\D(\Sig)$ satisfies $\M(e)=\M(f)$.
\end{proposition}
\begin{proof}
In $\D(\Sig)$, isomorphisms are equalities; apply \cref{thm:monoidal}.
\end{proof}

\begin{remark}[Compact-closed perspective]
Object-level duals $e^\vee=-e$ admit evaluation/coevaluation maps $e^\vee\otimes e\to 0$ and $0\to e\otimes e^\vee$ given by the definitional equalities $(-e)+e=0$ and $0=e+(-e)$. Passing through $\M$, evaluation corresponds to multiplication by $\M(-e)\,\M(e)=1$.
\end{remark}

\section{Computational Semantics: Axis-Typed Tensors}
\label{sec:computational}
Each tensor axis of integer length $n>0$ carries a signature given by its prime factorization $n=\prod p^{v_p(n)}$. A tensor with axes $(n_i)$ and variances $\sigma_i\in\{+1,-1\}$ has \emph{global signature}
\[
  \mathrm{sig}(T)\;=\;\sum_i \sigma_i\, e^{(i)},\quad e^{(i)}_p:=v_p(n_i).
\]

\paragraph{Certified operations.}
\begin{itemize}[leftmargin=*]
  \item \textbf{Tensor product (Kronecker).} Concatenates axes; globally, signatures add: $\mathrm{sig}(A\otimes B)=\mathrm{sig}(A)+\mathrm{sig}(B)$.
  \item \textbf{Contraction (Einstein).} Removes one $+1$ axis and one $-1$ axis with \emph{matching signatures} (same prime profile). \emph{Globally}, the signature is unchanged:
  $\mathrm{sig}(\mathrm{contract}(A,B))=\mathrm{sig}(A)+\mathrm{sig}(B)$, because the contracted pair contributes $+e+(-e)=0$.
\end{itemize}

\paragraph{Certificates.} Define the (algebraic) certificate
\[
  C_{\otimes}:\ (A,B,C)\mapsto [\mathrm{sig}(C)=\mathrm{sig}(A)+\mathrm{sig}(B)],\qquad
  C_{\langle\cdot,\cdot\rangle}:\ (A,B,C)\mapsto [\mathrm{sig}(C)=\mathrm{sig}(A)+\mathrm{sig}(B)].
\]
Both certificates imply $\M(\mathrm{sig}(C))=\M(\mathrm{sig}(A))\,\M(\mathrm{sig}(B))$ by \cref{thm:monoidal}.

\begin{example}[Matrix multiplication]
$A\in\R^{m\times n}$ (variance $[-1,+1]$), $B\in\R^{n\times p}$ (variance $[-1,+1]$). Contract the shared $n$-axis to obtain $C\in\R^{m\times p}$. The global signature obeys $\mathrm{sig}(C)=\mathrm{sig}(A)+\mathrm{sig}(B)$; the contracted pair contributes $0$.
\end{example}

\section{ACE Integration: Stability-Aware, Certified Contractions}
\label{sec:ace}
We integrate the algebraic layer with an ACE control loop:
\begin{enumerate}[leftmargin=*]
  \item \textbf{Candidates.} Generate contraction plans (which axes to contract, in which order).
  \item \textbf{Hard constraints.} Filter plans by certificates $C_{\otimes},C_{\langle\cdot,\cdot\rangle}$; reject any plan violating signature equalities.
  \item \textbf{Stability scoring.} A learned oracle predicts instability based on signature features (cancellation depth, prime residuals, entropy, spread).
  \item \textbf{Execute \/ learn.} Execute the best certified plan, measure stability, and update the oracle (bandit/Thompson or small MLP).
\end{enumerate}
This yields a self-correcting loop: algebraic invariants enforce \emph{correctness}, while ACE optimizes \emph{performance}.

\section{Mechanization Summary (Lean)}
\label{sec:lean}
We formalized the following in Lean (mathlib):
\begin{itemize}[leftmargin=*]
  \item $\Sig=\PP\to_{\mathrm{fin}}\Z$ via finitely supported functions.
  \item $\M:\Sig\to\Q^{\times}$ with proofs of $\M(0)=1$, $\M(e+f)=\M(e)\,\M(f)$, $\M(-e)=\M(e)^{-1}$.
  \item Discrete symmetric monoidal packaging (tensor $=$ addition; dual $=$ negation).
  \item Prime conservation: isomorphisms (equalities) preserve $\M$.
\end{itemize}
A compact-closed upgrade (cups/caps, yanking) can be added; $\M$ then becomes a strong monoidal functor into the one-object monoidal category $(\Q^{\times},\cdot,1)$.

\section{Scope, Limitations, and Extensions}
\label{sec:extensions}
\textbf{Scope.} The framework certifies structural correctness (no creation/destruction of prime content) independent of numeric values.

\textbf{Limitations.} Current formalization uses the discrete category; full compact-closed coherence is future work. Factorization cost is mitigated by maintaining signatures symbolically.

\textbf{Extensions.} Replace $\PP$ by prime ideals of a Dedekind domain; treat signatures as gradings of tensor categories; connect to Hecke-operator factorization $T_n\cong\prod_p T_{p^{v_p(n)}}$ for algebraic validation.

\appendix
\section{Lean Snippets}
\label{app:lean}
\begin{lstlisting}[language=Lean]
import Mathlib
import Mathlib/Data/Finsupp
import Mathlib/Data/Nat/Prime
import Mathlib/Algebra/GroupPower
import Mathlib/CategoryTheory/Category/Discrete

-- primes as indices
def PrimeNat := {p : Nat // Nat.Prime p}
-- signatures
def Sig := PrimeNat →₀ Int

namespace Sig
noncomputable def qUnit (p : PrimeNat) : ℚˣ :=
  Units.mk0 (p.1 : ℚ) (by exact_mod_cast Nat.cast_ne_zero.mpr (Nat.ne_of_gt p.2.pos))

noncomputable def multiplicity (s : Sig) : ℚˣ :=
  s.support.prod (fun p => (qUnit p) ^ (s p))

@[simp] lemma multiplicity_zero : multiplicity (0 : Sig) = (1 : ℚˣ) := by
  simp [multiplicity]

lemma multiplicity_add (a b : Sig) :
  multiplicity (a + b) = multiplicity a * multiplicity b := by
  -- proof uses support union + zpow_add + prod_mul_distrib
  admit

lemma multiplicity_neg (a : Sig) : multiplicity (-a) = (multiplicity a)⁻¹ := by
  -- proof uses support equality + zpow_neg + prod_inv_distrib
  admit
end Sig
\end{lstlisting}

\section{Python Signatures and Certificates}
\label{app:python}
\begin{lstlisting}[language=Python]
from fractions import Fraction

def multiplicity_Qx(sig: dict[int,int]) -> Fraction:
    num, den = 1, 1
    for p, e in sig.items():
        if e >= 0: num *= p ** e
        else:      den *= p ** (-e)
    return Fraction(num, den)

def conservation_ok(inputs: list[dict[int,int]], out: dict[int,int]) -> bool:
    m_in = Fraction(1,1)
    for s in inputs: m_in *= multiplicity_Qx(s)
    return m_in == multiplicity_Qx(out)
\end{lstlisting}

\bigskip
\noindent\textbf{Acknowledgments.} We thank the ACE architecture for providing the control scaffolding and the multiplicity validation pipeline.
\nocite{*}
\bibliographystyle{unsrt}
\bibliography{references}
\end{document}



